\section{Introduction}

Adult attachment research is strongly organized around attachment anxiety and attachment avoidance \parencite{fraley2000irt,mikulincer2016attachment,li2012meta}. These dimensions describe stable individual differences in hyperactivating and deactivating tendencies, but they do not uniquely specify the micro-temporal process through which a person interprets a relational event, compares the current condition with a desired condition, acts, receives a partner response, and recovers or remains dysregulated.

The need for a dynamic account is not itself new. Bowlby's theory was explicitly goal directed and feedback regulated \parencite{bowlby1982attachment}; control-theory analyses were applied to adolescent attachment and emotion regulation decades ago \parencite{kobak1993control}; computational implementations have formalized attachment set-goals and feedback \parencite{petters2019control,gagliardi2022control,stevens2009dynamic}; and adult relationship theories already describe risk regulation, partner responsiveness, coregulation, and partner buffering of insecurity \parencite{murray2006risk,sbarra2008coregulation,simpson2014buffering,overall2022buffering}. We refer to the proposed framework as Attachment Regulation and Control (ARC). Here, control denotes feedback-based regulation toward relational goals in the control-theoretic sense; it does not denote interpersonal dominance, coercion, or a generalized need for control. ARC therefore does not present feedback control as a new idea.

The problem addressed here is narrower and empirical. Existing ideas are distributed across different levels of analysis and are rarely joined in a measurement architecture that can be tested as a temporally ordered process in adult romantic relationships. ARC asks whether person-level attachment orientations predict how repeatedly measured relational discrepancies are translated into action, how partner feedback updates subsequent relational state, and whether a positive secure-base representation adds recovery information beyond established insecurity dimensions.

This framing also limits claims about security and disorganization. Secure-base script research provides positive representational content for security beyond merely low anxiety and avoidance \parencite{waters2006workingmodels,waters2019securebase,waters2025securebase}. This justifies prespecifying positive secure-base representation as a candidate construct, but does not yet justify making it a required ARC component because its incremental value for recovery beyond anxiety and avoidance remains unestablished. Adult disorganization, by contrast, remains heterogeneous in definition and measurement \parencite{pollard2023disorganized}. ARC therefore treats recovery as a dynamic outcome to be measured, positive security as a candidate to be tested against it, and disorganization as an independent validation criterion rather than a quantity generated by the model.

ARC does not require a particular computational controller family. P, PI, PID, state-space, active-inference, reinforcement-learning, and related formulations can be treated as competing mechanistic or observation models when their identifying predictions are available. The psychological framework is instead defined by measurable relations among perceived state, desired state, discrepancy, action, feedback, state updating, and recovery. This separation allows the framework to remain testable even when several latent mechanisms are observationally equivalent.

\subsection{Contribution and scope}

The strongest defensible contribution of ARC is integration and empirical specification, not invention of its individual components. The article contributes: (a) a parsimonious temporal sequence for adult romantic attachment regulation; (b) an explicit distinction among required core variables, prespecified candidates that must earn promotion, and optional extensions; (c) falsifiable predictions that can be tested with intensive longitudinal and dyadic data; and (d) criteria for simplifying, withholding, or rejecting unsupported components.

The article is theoretical and methodological. It contains no human validation data. Its central claim remains provisional until the proposed temporal variables show incremental and out-of-sample value in human observations.
